\newpage
\phantomsection
\label{prf:two-element-system-not-ordered-field}
\LRAProofFor{thm:two-element-system-not-ordered-field}

\begin{remark*}[Return]
\hyperref[thm:two-element-system-not-ordered-field]{Return to Theorem}
\end{remark*}

\begin{theorem*}[The Two-Element System Is Not an Ordered Field]
There is no order on \(\mathbb{B}_2\) that makes \(\mathbb{F}_2\) an ordered
field.
\end{theorem*}

\begin{proof}[Professional Standard Proof]
\LRAProofBodyStart
\textbf{Professional Standard Proof.}
In any ordered field, \(1>0\), and addition preserves strict order. Adding
\(1\) to both sides gives \(1+1>0+1\). In \(\mathbb{F}_2\), however,
\(1+1=0\) and \(0+1=1\), so this would imply \(0>1\), contradicting
\(1>0\). Thus no compatible field order exists.
\end{proof}

\begin{proof}[Detailed Learning Proof]
\LRAProofBodyStart
\textbf{Detailed Learning Proof.}
Assume an order makes \(\mathbb{F}_2\) an ordered field. Ordered fields have
positive square \(1=1^2>0\), so \(1>0\). The order must be compatible with
addition, so adding \(1\) to the inequality \(1>0\) gives
\[
1+_2 1 > 0+_2 1.
\]
The two-element addition table gives \(1+_2 1=0\) and \(0+_2 1=1\). Hence
\(0>1\), contradicting \(1>0\). Therefore no such order exists.
\end{proof}

\begin{remark*}[Proof structure]
The obstruction is algebraic: the identity \(1+1=0\) is incompatible with
ordered-field positivity and additive order preservation.
\end{remark*}

\begin{dependencies}
\begin{itemize}
  \item \hyperref[thm:two-element-system-is-field]{The Two-Element System Is a Field}
  \item \hyperref[def:abstract-ordered-field]{Ordered Field}
\end{itemize}
\end{dependencies}

\clearpage
